%=======================================================================
%  The Experiment in Depth
%=======================================================================
\section{The Experiment in Depth}
\label{sec:experiment}

This section describes exactly what was run. Section~\ref{sec:exp-anatomy}
traces one real trial end to end with its actual numbers, and
Sections~\ref{sec:exp-one}--\ref{sec:exp-sensitivity} give the protocol
for each of the five experiments.

\subsection{The five experiments}

\begin{table}[htbp]
\centering
\caption{The five experiments, what each isolates, and what it produces.}
\label{tab:experiments}
\small
\begin{tabular}{@{}p{0.5cm}p{3.0cm}p{5.0cm}p{4.6cm}@{}}
\toprule
\textbf{\#} & \textbf{Experiment} & \textbf{Isolates} & \textbf{Produces} \\
\midrule
E1 & Paired randomised trials &
Countermeasure effect, holding the scenario fixed &
\TotalTrials{} matched pairs per vector; RQ1--RQ3 \\
\addlinespace
E2 & Concrete mined elections &
Whether abstract probabilities correspond to real winner flips &
\ElectionTrials{} mined elections per vector; ground truth \\
\addlinespace
E3 & Benchmark scenarios &
Behaviour at five interpretable, fixed operating points &
Quotable per-scenario numbers; a negative control \\
\addlinespace
E4 & Threshold solving &
The parameter value at which each attack stops being viable &
Seven thresholds per scenario; RQ4 \\
\addlinespace
E5 & Coalition and sensitivity &
Whether combining vectors beats the best single vector, and whether
that survives the cost assumptions &
231-point simplex search; \SensitivityConfigs-configuration sweep; RQ5 \\
\bottomrule
\end{tabular}
\end{table}

\subsection{Anatomy of a single trial}
\label{sec:exp-anatomy}

The clearest way to describe the experiment is to follow one trial
through it. Trial \WEtrial{} is used below. It is an ordinary trial
drawn by the sampler, not a constructed illustration, and it was chosen
because it happens to exhibit all three of the behaviours the model is
built to distinguish.

\subsubsection*{Step 1 --- Draw the scenario}

The generator draws one scenario from the ranges of
\tabref{tab:parameters}, using the trial's own stream from the seeded
generator:

\begin{center}
\small
\begin{tabular}{@{}llll@{}}
\toprule
$N = \WEN$ voters & $m = \WEm$ votes & $B = \WEB$ & $R = \WER$ \\
$P = \WEP$ & $S = \WES$ identities & $A = \WEA$ ballots & \\
\bottomrule
\end{tabular}
\end{center}

The manipulation threshold is $m = \WEm$: the attack succeeds if and
only if at least that many fraudulent votes survive every
countermeasure and reach the count. Every number above is now
\emph{fixed} for both arms of this pair. Only the acceptance and
detection probabilities will differ.

\subsubsection*{Step 2 --- Bribery: the feasibility constraint}

The attacker can approach at most
\[
T = \min\!\left(N, \left\lfloor \frac{B}{R} \right\rfloor\right)
  = \min\!\left(\WEN,\ \left\lfloor \frac{\WEB}{\WER} \right\rfloor\right)
  = \WET \text{ voters},
\]
so the budget binds well before the electorate does. Note already that
$T = \WET$ exceeds $m = \WEm$ --- the attack is \emph{arithmetically}
possible --- but only just, which is what makes the survival
probability decisive.

\subsubsection*{Step 3 --- Bribery: the rational-voter gate}

This is where the two arms diverge, and it is the most consequential
step in the model.

\emph{Baseline arm}, $q = \WEqBase$:
\[
\Uv = (1-q)R - qP
    = (1 - \WEqBase)(\WER) - (\WEqBase)(\WEP)
    = \WEUvBase > 0
    \;\Rightarrow\; \text{voter } \WEdecisionBase\text{s}.
\]

\emph{Protected arm}, $q = \WEqProt$:
\[
\Uv = (1 - \WEqProt)(\WER) - (\WEqProt)(\WEP)
    = \WEUvProt < 0
    \;\Rightarrow\; \text{voter } \WEdecisionProt\text{s}.
\]

The attacker's budget, the electorate and the margin are identical in
both arms. The bribe is simply no longer worth accepting. Under the
protected configuration the attack contributes zero fraudulent votes
not because they were detected and removed, but because they were never
cast.

\subsubsection*{Step 4 --- Bribery: distribution and success probability}

In the baseline arm each bribed vote survives with probability
$1 - q = \WEpsurviveBase$, so
\[
X \sim \Bin(\WET,\ \WEpsurviveBase),
\qquad
\mu = \Ex[X] = \WEmuBase .
\]

The mean $\mu = \WEmuBase$ sits \emph{below} the threshold $m = \WEm$,
so the attack succeeds only on an upper-tail excursion. The exact
probability is
\[
\Prob(X \geq \WEm) = \WEexactBase .
\]

Three methods are now applied to this same quantity:

\begin{center}
\small
\begin{tabular}{@{}lll@{}}
\toprule
\textbf{Method} & \textbf{Value} & \textbf{Role} \\
\midrule
Exact binomial survival function & \WEexactBase & ground truth \\
Monte Carlo, \MonteCarloTrials{} draws & \WEmcBase & independent numerical check \\
Chernoff upper bound (M7) & \WEchernoffBase & provable ceiling \\
\bottomrule
\end{tabular}
\end{center}

All three orderings that the mathematics requires hold: the Monte Carlo
estimate is close to the exact value, and the bound lies above it. The
bound is defined here precisely because $m > \mu$. In the protected arm
the gate is closed, so $\Prob = \WEexactProt$ and the bound is
\WEchernoffProt{} --- correctly reported as absent rather than as zero,
since no bound was computed.

\subsubsection*{Step 5 --- Sybil: two gates in series}

The same trial, same threshold $m = \WEm$, with $S = \WES$ forged
identities attempted. A forged identity must pass registration
\emph{and} escape detection to place a counted vote:
\[
W \sim \Bin\big(S,\; q_s(1-p_d)\big).
\]

\emph{Baseline}, $q_s = \WEqsBase$, $p_d = \WEpdBase$:
$\Ex[W] = \WEsybilmuBase$, which is more than twice the threshold, so
$\Prob(W \geq m) = \WEsybilPBase$ --- effectively certain.

\emph{Protected}, $q_s = \WEqsProt$, $p_d = \WEpdProt$:
$\Ex[W] = \WEsybilmuProt$, now well below the threshold, so
$\Prob(W \geq m) = \WEsybilPProt$.

The mechanism differs from bribery's. Nothing refuses; the two gates
simply tighten, the mean crosses from one side of the threshold to the
other, and the probability follows it sharply because the binomial
concentrates around its mean.

\subsubsection*{Step 6 --- Ballot stuffing: structural infeasibility}

The same trial attempts $A = \WEA$ fraudulent ballots against a
threshold of $m = \WEm$. Since $A < m$, the attack cannot reach the
threshold even if every single ballot were accepted and undetected:
\[
\Prob(W \geq \WEm) = \WEballotPBase \text{ (baseline)},
\qquad
\WEballotPProt \text{ (protected)} .
\]

The expected surviving counts --- \WEballotmuBase{} and
\WEballotmuProt{} --- confirm how far short the attack falls. This is
not a detection failure but a capacity limit, and it is the phenomenon
that Section~\ref{sec:coalition} shows to be escapable by combining
vectors.

\subsubsection*{What this one trial demonstrates}

\begin{finding}
\finhead{Three distinct failure mechanisms, in one trial}
Trial \WEtrial{} shows the model distinguishing three different reasons
an attack fails. Bribery fails in the protected arm because the bribe
stops being \emph{worth accepting} ($\Uv = \WEUvProt$). The Sybil
attack fails because the expected surviving count crosses below the
threshold (\WEsybilmuBase{} $\rightarrow$ \WEsybilmuProt{} against
$m = \WEm$). Ballot stuffing fails in both arms because fewer ballots
are attempted than the threshold requires ($A = \WEA < m = \WEm$) ---
a structural limit no detection setting affects. A model reporting only
aggregate success rates would collapse these three into one number and
lose the distinction that tells a defender what to fix.
\end{finding}

\subsection{Experiment 1: paired randomised trials}
\label{sec:exp-one}

\textbf{Procedure.} Repeat the trial above \TotalTrials{} times per
vector, drawing a fresh scenario each time and evaluating both arms.
Per trial, record: the drawn parameters; the rational-voter decision;
$T$; the exact, Monte Carlo and Chernoff values; expected surviving
counts; detection probabilities; and attacker cost.

\textbf{Output.} 5{,}000 rows per vector (baseline and protected for
each of \TotalTrials{} pairs), 15{,}000 rows in the combined
Chernoff-validation table, and an aggregate summary.

\textbf{Analysis.} Paired differences are aggregated into means with
bootstrap confidence intervals (10{,}000 resamples), and the exact
value is compared against the Monte Carlo estimate and the Chernoff
bound trial by trial.

\textbf{What it controls for.} Because both arms share the scenario, any
difference is attributable to the configuration alone. What it does
\emph{not} control for is the choice of parameter ranges, which is why
the ranges are broad and why the benchmark scenarios of E3 exist as a
complementary fixed-point analysis.

\subsection{Experiment 2: concrete mined elections}
\label{sec:exp-two}

E1 produces probabilities. E2 produces winners.

\textbf{Procedure.} For each of \ElectionTrials{} trials per vector:

\begin{enumerate}[leftmargin=1.7em,itemsep=2pt]
\item Draw an electorate of 150--900 voters and a support level for
      candidate B in $[0.35, 0.65]$, biased toward close races so that
      the manipulation threshold is small enough for an attack to be
      interesting.
\item Every legitimate voter casts a ballot. The batch is sealed into a
      genuinely mined block: proof-of-work at difficulty 3, meaning the
      block hash must begin with three zeros, found by incrementing a
      nonce until it does.
\item Read the pre-attack tally \emph{from the chain} via
      \code{count\_votes()} --- not from a Python-side running total.
      This matters: the tally under test is the ledger's own.
\item Compute the flip threshold $m$ from the actual margin
      (\eqnref{eq:threshold}) and identify the losing candidate as the
      attacker's target.
\item Simulate each fraudulent attempt individually through the real
      pipeline: an acceptance gate (bribe accepted / registration
      passed / ballot validated), then a detection gate. Only votes
      that are accepted \emph{and} undetected are mined into the chain.
      Rejected and detected attempts are logged but never reach the
      ledger, which is what a real countermeasure does --- a rejected
      transaction is not a recorded vote.
\item Read the post-attack tally from the chain, determine the winner,
      and record whether it changed.
\item Run \code{is\_chain\_valid()} to confirm the ledger is still
      internally consistent: every hash correct, every link intact,
      every block satisfying the difficulty target.
\end{enumerate}

Both arms run against \emph{the same starting election}, so the pairing
of E1 is preserved: the two configurations are handed an identical
mined chain and an identical margin.

\textbf{What it establishes.} Two things. First, that the abstract
probabilities correspond to real outcomes rather than to an artefact of
the abstraction. Second --- and this is the finding the dissertation
turns on --- that chain integrity and election correctness are
independent: the validity check passes in every trial, including the
ones where the winner changed.

\subsection{Experiment 3: benchmark scenarios}
\label{sec:exp-three}

E1 averages over a distribution of scenarios, which is statistically
powerful and hard to reason about concretely. E3 fixes five
interpretable operating points --- a university election ($N = \SOneN$),
an organisation election ($N = \STwoN$), a national election
($N = \SThreeN$), a maximum-capability high-budget attack, and a
deliberately under-resourced low-budget attack --- and evaluates each
exactly, by Monte Carlo, by Chernoff bound and by the game-theoretic
optimiser.

S5 is a negative control. Its threshold is set far beyond what its
budget can reach, and the model is expected to report it as failing; a
model that reported S5 succeeding would be wrong in a way that no
amount of aggregate statistics would reveal. It reports
$P = \SFiveExactFlip$.

\subsection{Experiment 4: threshold solving}
\label{sec:exp-four}

For each benchmark scenario, the seven thresholds of
Section~\ref{sec:thresholds} are solved.

\textbf{Closed-form thresholds} (T1, T2) are evaluated directly. T2 is
also verified against the numerically computed zero contour of the
utility surface, and the two coincide (\figref{fig:utility}).

\textbf{Integer thresholds} (T3, T6, T7) are found by binary search
over a monotone function, which is exact. For T6 the search runs over
margins above the mean, where the Chernoff bound decreases
monotonically.

\textbf{Continuous thresholds} (T4, T5) are found by bisection on the
signed utility of attacking. Each solver first checks that the endpoints
bracket a crossing, and returns an explicit absence if they do not ---
which is what happens for $B^*$ under the protected configuration,
where no budget makes the attack profitable because the bribe is not
incentive-compatible at any budget.

Two supporting sweeps generate the data behind \figref{fig:critical}
and the break-even budget: utility against detection probability across
$q \in [0, 1]$, and utility against budget up to three times the
scenario budget, both for every scenario and both arms.

\subsection{Experiment 5: coalition and sensitivity}
\label{sec:exp-sensitivity}

\textbf{Coalition search.} For each scenario and arm, the three vectors
are placed on a common budget line by assigning unit costs: a bribe
costs $R$; a forged identity costs \IdentityCostRatio\% of $R$; an
injected ballot costs \BallotCostRatio\% of $R$. Structural ceilings
apply --- bribery at $N$ by \eqnref{eq:m3}, forged registrations at
\SybilCeilingPct\% of $N$, injected ballots at \BallotCeilingPct\% of
$N$. The budget simplex is then evaluated exhaustively at a 0.05 step,
231 allocations, each scored by the exact convolution of the three
component distributions. The best allocation is compared against the
best pure single-vector strategy, which is a corner of the same
simplex.

\textbf{Sensitivity sweep.} The four constants above have no empirical
basis, so the entire coalition analysis is re-run across their product:
identity cost in $\{0.2, 0.4, 0.8\}$, ballot cost in $\{0.05, 0.15,
0.40\}$, registration ceiling in $\{0.10, 0.20\}$ and ballot ceiling in
$\{0.15, 0.30\}$, for three scenarios --- \SensitivityConfigs{}
configurations in total. For each, the analysis records whether the
coalition still beats the best pure strategy.

This is the check that decides whether the coalition finding is a
result or an artefact of convenient parameters, and it is reported as
such in Section~\ref{sec:findings-coalition}.
